\documentclass[11pt]{article}
% Shared print layout. The web's provisional design and MathJax are independent.
\usepackage[utf8]{inputenc}
\usepackage[T1]{fontenc}
\usepackage[spanish,es-nodecimaldot]{babel}
\usepackage{lmodern}
\usepackage{amsmath,amssymb}
\usepackage{graphicx}
\usepackage{geometry}
\geometry{a4paper,left=20mm,right=20mm,top=25mm,bottom=22mm,headheight=16pt,headsep=8mm}
\usepackage{fancyhdr,lastpage}
\usepackage{enumitem}
\usepackage{xcolor}
\usepackage{pgfplots}
\pgfplotsset{compat=1.18}
\usepgfplotslibrary{fillbetween}
\usepackage[hidelinks,pdfencoding=auto]{hyperref}
\hypersetup{pdflang=es-ES,pdfauthor={},pdfsubject={Matemáticas Aplicadas a las Ciencias Sociales II}}
\pagestyle{fancy}
\fancyhf{}
\fancyhead[L]{\small Matemáticas Aplicadas a las Ciencias Sociales II}
\fancyhead[R]{\small 2.\textsuperscript{o} Bachillerato}
\fancyfoot[L]{\footnotesize Subida de nota de curso | Edición revisada 1.0}
\fancyfoot[R]{\footnotesize Página \thepage\ de \pageref{LastPage}}
\renewcommand{\headrulewidth}{0.4pt}
\setlength{\parindent}{0pt}
\setlength{\parskip}{0.5em}
\setlist{itemsep=0.3em,topsep=0.3em}
\newcommand{\apartado}[1]{\subsection*{#1}}
\pgfplotsset{
  revision/.style={
    width=0.94\linewidth,height=7.5cm,
    axis lines=middle,grid=major,grid style={gray!18},
    tick label style={font=\small},label style={font=\small},
    xlabel={$x$},ylabel={$y$},samples=160,
    legend style={at={(0.5,-0.16)},anchor=north,draw=none,font=\small},
    clip=true
  }
}

\hypersetup{pdftitle={Subida de nota de curso - Soluciones revisadas}}
\begin{document}

\section*{Bloque A: continuidad, monotonía y asíntotas}
Soluciones originales revisadas. Se conservan cuatro ejercicios, 90 minutos y 10 puntos; 0,25 de expresión incluidos por ejercicio y condiciones de calculadora del original.
\begin{enumerate}[label=(\alph*)]
\item En 2, el valor y el límite derecho son 1; el izquierdo vale $a-4$. Continuidad exige $a-4=1$, luego $\boxed{a=5}$.

Para ese valor, las derivadas laterales son 0 y $-1$; por tanto, no es derivable en 2. El cociente incremental izquierdo es $h\to0$; el derecho es $-1/(1+h)\to-1$.
\item Ahora se usa $a=4$, como pide este apartado: la rama izquierda es $(x-2)^2$ para $x<2$, y la derecha es $1/(x-1)$ solo para $x\geq2$.

\textbf{Corrección M01:} el dominio de la función a trozos es $\mathbb R$, pues en torno a 1 rige el polinomio, no la expresión racional. En particular $f(1)=1$ y $\lim_{x\to1}f(x)=1$: \textbf{no hay asíntota vertical en 1} ni en otro punto.

Para $x<2$, $f'(x)=2x-4<0$; para $x>2$, $f'(x)=-1/(x-1)^2<0$. Es estrictamente decreciente en cada intervalo $(-\infty,2)$ y $[2,\infty)$ por separado. No es decreciente en todo $\mathbb R$: $f(3/2)=1/4<f(2)=1$ aunque $3/2<2$.

En 2 hay salto desde límite izquierdo 0 hasta valor y límite derecho 1. No es una discontinuidad infinita. El salto invalida la conclusión global del original pese a las derivadas negativas en las ramas.

Al tender $x$ a $+\infty$, $f(x)\to0$: asíntota horizontal $\boxed{y=0}$ por la derecha. Por la izquierda la rama es cuadrática, sin asíntota horizontal ni oblicua lineal; no se extiende la rama racional fuera de su intervalo.
\end{enumerate}
\clearpage
\section*{Bloque B: confitería y beneficio máximo}
Sean $x$ las tartas de chocolate e $y$ las de chocolate y merengue. Se conserva el modelo original:
\[
\max B=10x+12y,\qquad
\begin{cases}
x+2y\leq100,\\
2x+y\leq80,\\
y\leq46,\\
x,y\geq0.
\end{cases}
\]
\textbf{Corrección M37:} deben comprobarse todas las restricciones antes de comparar beneficios. El recinto tiene los cinco vértices siguientes:
\[
\begin{array}{c|r}
(x,y)&B\text{ (euros)}\\\hline
(0,0)&0\\(40,0)&400\\(20,40)&680\\(8,46)&632\\(0,46)&552
\end{array}
\]
\begin{center}\begin{tikzpicture}\begin{axis}[revision,height=5cm,xmin=0,xmax=55,ymin=0,ymax=55]
\addplot[fill=blue!15,draw=blue,thick] coordinates {(0,0)(40,0)(20,40)(8,46)(0,46)(0,0)};
\addplot[domain=0:55,dashed]{(100-x)/2};\addplot[domain=12.5:40,dotted]{80-2*x};\addplot[domain=0:55,dashdotted]{46};
\addplot[only marks,blue] coordinates {(20,40)};
\end{axis}\end{tikzpicture}\end{center}
Máximo único $\boxed{680\text{ euros en }(20,40)}$: 20 tartas de chocolate y 40 de chocolate y merengue. Consumen 100 kg de bizcocho, 80 de crema y 40 de merengue, dentro de los 46 disponibles.

El óptimo original $(17,46)$ consume 109 kg de bizcocho, de modo que no es factible. $(50,0)$ incumple la crema; $(0,40)$ es factible pero no es vértice del recinto. El modelo continuo alcanza un óptimo entero, por lo que también resuelve el problema si se exigen tartas enteras.
\clearpage
\section*{Bloque C: incompatibilidad y probabilidades}
\begin{enumerate}[label=(\alph*)]
\item Si $P(A)>0$ y $P(B)>0$, independencia exigiría $P(A\cap B)=P(A)P(B)>0$. La incompatibilidad exige intersección vacía y, por tanto, probabilidad 0. No pueden darse ambas condiciones simultáneamente bajo esas hipótesis. El requisito de probabilidades no nulas es esencial.
\item Se interpreta $P(A/B)$ como $P(A\mid B)$, conforme a la solución original. Dado que $P(B)=0,5>0$,
\[
P(A\cap B)=0,2\cdot0,5=\boxed{0,1}.
\]
\[
P(A\cup B)=0,3+0,5-0,1=\boxed{0,7}.
\]
Por De Morgan,
\[
P(A^c\cup B^c)=1-P(A\cap B)=\boxed{0,9}.
\]
\[
P(A-B)=P(A\cap B^c)=0,3-0,1=\boxed{0,2}.
\]
Como comprobación, los cuatro átomos $A\cap B$, $A\cap B^c$, $A^c\cap B$, $A^c\cap B^c$ tienen probabilidades $0,1;0,2;0,4;0,3$, todas no negativas y con suma 1.
\end{enumerate}
\clearpage
\section*{Bloque D: desfase del autobús}
Se conserva la varianza poblacional 4, luego $\sigma=2$ minutos. Se supone muestreo aleatorio con independencia acorde al modelo normal. La normalidad de la media es exacta bajo ese modelo aunque $n=10$; no se justifica por una regla universal de tamaño grande.
\begin{enumerate}[label=(\alph*)]
\item \textbf{Corrección M38:} los diez datos, incluido $-0,2$, suman 30,9 minutos, no 31,9:
\[
\bar x=30,9/10=\boxed{3,09\text{ minutos}}.
\]
Para confianza del 97\%, $z_c=\Phi^{-1}(0,985)\approx2,17009038$:
\[
IC=\bar x\pm z_c\frac2{\sqrt{10}}
=\boxed{(1,71751434;\ 4,46248566)}\text{ minutos}.
\]
Con tabla $z_c\approx2,17$ el margen es 1,37242850 y, a dos decimales, el intervalo es $(1,72;4,46)$ minutos. No se centra en la media errónea 3,19 del original.
\item 30 segundos son 0,5 minutos. Se pide margen estrictamente inferior a esa cantidad. Con confianza del 95\%,
\[
q=\left(\frac{\Phi^{-1}(0,975)\,2}{0,5}\right)^2\approx61,46334113,
\qquad n>q.
\]
\[
\boxed{n_{\min}=\lfloor q\rfloor+1=62}.
\]
Con 61 observaciones el margen es aproximadamente 0,50190 minutos, y con 62, 0,49783; solo el segundo cumple. El cuantil de tabla 1,96 conduce al mismo entero.

Aumentar la confianza, manteniendo margen y desviación, aumenta la cota real q. El tamaño mínimo entero no disminuye; puede mantenerse igual para aumentos pequeños de confianza por el redondeo a enteros. Se precisa así el aumento afirmado en el original, sin exigir crecimiento estricto en cada incremento de confianza.
\end{enumerate}
Los intervalos tienen la interpretación de cobertura del procedimiento bajo muestreo repetido; no asignan probabilidad posterior a una media fija. No se inventan tamaño de población ni correcciones finitas sin datos.
\end{document}
